\section{Structure sheaf as a sheaf of abelian groups, sheafification and Ext}
\label{mk-chap:Cohomology_StructureSheafAb}

Let \(X\) be a topological space. Abelian-group-valued presheaves on the topology of opens admit sheafification, and their sheaf category admits \(\Ext\). Together with \ref{mk-chap:Cohomology_SheafCompose}, these facts allow the structure sheaf of a scheme to be regarded as a sheaf of abelian groups.

\begin{enumerate}
  \item Sheafification for abelian-group-valued presheaves on opens.
  \item \(\Ext\) in the resulting sheaf category.
  \item The underlying abelian-group-valued structure sheaf \(\mathcal O_C\).
\end{enumerate}

\section{Sheafification of abelian-group-valued sheaves on opens}

\begin{theorem}[Sheafification on the topology of opens, abelian-group valued]\label{mk-thm:HasSheafify_Opens_AddCommGrp}
  \lean{AlgebraicGeometry.Cohomology.instHasSheafify_Opens_AddCommGrp}
  \leanok
  \dcref{custom}

  Let \(X\) be a topological space. Then the Grothendieck topology of opens of \(X\) admits sheafification for sheaves valued in the category of abelian groups.
\end{theorem}

\begin{proof}
  \leanok


The category of abelian groups is a Grothendieck abelian category. Hence abelian-group-valued presheaves on the site of opens admit the usual sheafification, obtained from the plus construction, and the sheafification functor is left exact.
\end{proof}

\section{\(\Ext\) for sheaves on opens}

\begin{theorem}[\(\Ext\) on the topology of opens, abelian-group valued]\label{mk-thm:HasExt_Sheaf_Opens_AddCommGrp}
  \lean{AlgebraicGeometry.Cohomology.instHasExt_Sheaf_Opens_AddCommGrp}
  \leanok
  \dcref{custom}

  \uses{mk-thm:HasSheafify_Opens_AddCommGrp}

  Let \(X\) be a topological space. Then the category of sheaves of abelian groups on the topology of opens of \(X\) admits \(\Ext\) groups.
\end{theorem}

\begin{proof}
  \leanok


By \ref{mk-thm:HasSheafify_Opens_AddCommGrp}, sheaves of abelian groups form a Grothendieck abelian category. In particular, this category has enough injectives, so the right derived functors of \(\Hom\) define its \(\Ext\) groups.
\end{proof}

\section{Structure sheaf as a sheaf of abelian groups}

\begin{definition}[Structure sheaf as a sheaf of abelian groups]\label{mk-def:Scheme_toAbSheaf}
  \lean{AlgebraicGeometry.Scheme.toAbSheaf}
  \leanok
  \dcref{custom}

  \uses{mk-thm:HasSheafCompose_forget, mk-thm:HasSheafify_Opens_AddCommGrp,
        mk-thm:HasExt_Sheaf_Opens_AddCommGrp}

  Let \(C\) be a scheme. The structure sheaf of \(C\), viewed as a sheaf of abelian groups on the topology of opens of the underlying topological space of \(C\), is obtained by post-composing the commutative-ring-valued structure sheaf with the forgetful composite \(\mathrm{CommRing} \to \mathrm{Ring} \to \mathrm{Ab}\) of \ref{mk-chap:Cohomology_SheafCompose}.
\end{definition}

\section{Cohomology of the structure sheaf}

The preceding constructions define the abelian-group-valued cohomology of the structure sheaf:
\[
  H^n(C, \mathcal O_C) \;:=\; \text{toAbSheaf}(C).H\,n,
\]
where \(H\) denotes sheaf cohomology. Each \(H^n(C, \mathcal O_C)\) is therefore canonically an abelian group.
